% One-page key-equation sheet for the Lagrange points class (317 class 24, Fri Nov 6 2026).
% Printed for the students and imported into the GoodNotes chapter deck as the page Michael inks on.
% Equation numbers come from AstroLagrange.aux through xr-hyper, so build the handout first:
%   latexmk -pdf -silent AstroLagrange.tex && latexmk -pdf -silent AstroLagrange-keysheet.tex
\documentclass[10pt]{article}
%% Document setup for Syllabus and HW

\textwidth=7in
\textheight=9.5in
\topmargin=-1in
\headheight=0in
\headsep=.5in
\hoffset  -.85in

\usepackage{amsmath} % for \text{}
\usepackage{amsfonts} % For \text{}
\usepackage{amssymb} % For \text{}
\usepackage{amsthm}
\usepackage{bigints}

%% Stolen from Nathan Baker
\newcommand{\erf}{{\mathrm{erf}}}
\newcommand{\erfc}{{\mathrm{erfc}}}
\newcommand{\erfi}{{\mathrm{erfi}}}
\newcommand{\argh}{{\mathrm{arg}}}
\newcommand{\atan}{{\mathrm{atan}}}
\newcommand{\acos}{{\mathrm{acos}}}
\newcommand{\asin}{{\mathrm{asin}}}
\newcommand{\mat}[1]{\,\underline{\underline{#1}}\,}
\newcommand{\abs}[1]{\left| #1 \right|}
\newcommand{\norm}[1]{\left\| #1 \right\|}
\newcommand{\order}[1]{{\mathcal{O}} \left( #1 \right)}
\newcommand{\op}[1]{{\mathcal{#1}}}
\newcommand{\myfile}[1]{\texttt{#1}}
\newcommand{\myvar}[1]{\textsf{#1}}
\newcommand{\mean}[1]{{\left\langle {#1} \right\rangle}}

% The master boolean -- set this to "masterfalse" or "mastertrue"
%\newif\ifmaster \masterfalse
\newif\ifmaster \mastertrue

\newcommand{\hidesolution}[1]{ \ifmaster { {#1} } \else { {Go to class.} } \fi }
\newcommand{\hideanswer}[1]{ \ifmaster { #1 } \else {} \fi }
\newcommand{\solution}[1]{ \begin{proof}[Solution] \hidesolution{#1} \end{proof} }
\newcommand{\answer}[1]{ \hideanswer{\begin{proof}[Answer] #1 \end{proof}} }

%% MGL
\newcommand{\CC}{\mathbb{C}}
\newcommand{\N}{\mathbb{N}}
\newcommand{\Ham}{\mathcal{H}}
%\newcommand{\vect}[1]{\,\vec{\mathbf{#1}}\,}
\newcommand{\vect}[1]{\,\mathbf{#1}\,}
\usepackage{tabularx}
\newcommand{\vtau}{\vect{\tau}}
\newcommand{\vomega}{\vect{\omega}}
\newcommand{\va}{\vect{a}}
\newcommand{\vb}{\vect{b}}
\newcommand{\vc}{\vect{c}}
\newcommand{\vd}{\vect{d}}
\newcommand{\ve}{\vect{e}}
\newcommand{\vf}{\vect{f}}
\newcommand{\vg}{\vect{g}}
\newcommand{\vh}{\vect{h}}
\newcommand{\vi}{\vect{i}}
\newcommand{\vj}{\vect{j}}
\newcommand{\vk}{\vect{k}}
\newcommand{\vl}{\vect{l}}
\newcommand{\vm}{\vect{m}}
\newcommand{\vn}{\vect{n}}
\newcommand{\vo}{\vect{o}}
\newcommand{\vp}{\vect{p}}
\newcommand{\vq}{\vect{q}}
\newcommand{\vr}{\vect{r}}
\newcommand{\vs}{\vect{s}}
\newcommand{\vt}{\vect{t}}
\newcommand{\vu}{\vect{u}}
\newcommand{\vv}{\vect{v}}
\newcommand{\vw}{\vect{w}}
\newcommand{\vx}{\vect{x}}
\newcommand{\vy}{\vect{y}}
\newcommand{\vz}{\vect{z}}

\newcommand{\vA}{\vect{A}}
\newcommand{\vB}{\vect{B}}
\newcommand{\vC}{\vect{C}}
\newcommand{\vD}{\vect{D}}
\newcommand{\vE}{\vect{E}}
\newcommand{\vF}{\vect{F}}
\newcommand{\vG}{\vect{G}}
\newcommand{\vH}{\vect{H}}
\newcommand{\vI}{\vect{I}}
\newcommand{\vJ}{\vect{J}}
\newcommand{\vK}{\vect{K}}
\newcommand{\vL}{\vect{L}}
\newcommand{\vM}{\vect{M}}
\newcommand{\vN}{\vect{N}}
\newcommand{\vO}{\vect{O}}
\newcommand{\vP}{\vect{P}}
\newcommand{\vQ}{\vect{Q}}
\newcommand{\vR}{\vect{R}}
\newcommand{\vS}{\vect{S}}
\newcommand{\vT}{\vect{T}}
\newcommand{\vU}{\vect{U}}
\newcommand{\vV}{\vect{V}}
\newcommand{\vW}{\vect{W}}
\newcommand{\vX}{\vect{X}}
\newcommand{\vY}{\vect{Y}}
\newcommand{\vZ}{\vect{Z}}
%\newcommand{\pd}[2]{\frac{\partial#1}{\partial#2}}
%\newcommand{\dd}[2]{\frac{d#1}{d#2}}
\newcommand{\dinline}[2]{d #1/d #2} % for derivatives
\newcommand{\pdinline}[2]{\partial#1/\partial#2}

\newcommand{\boldcent}[1] {\begin{center}\textbf{ #1 }\end{center}}
% ***********************************************************
% ******************* PHYSICS HEADER ************************
% ***********************************************************
% From http://www.dfcd.net/articles/latex/latex.html
% Version 2
%\documentclass[11pt]{article} 
\usepackage{amsmath} % AMS Math Package
\usepackage{amsthm} % Theorem Formatting
\usepackage{amssymb}	% Math symbols such as \mathbb
\usepackage{graphicx} % Allows for eps images
\usepackage{multicol} % Allows for multiple columns
%\usepackage[dvips,letterpaper,margin=0.75in,bottom=0.5in]{geometry}
% % Sets margins and page size
%\pagestyle{empty} % Removes page numbers
%\makeatletter % Need for anything that contains an @ command 
%\renewcommand{\maketitle} % Redefine maketitle to conserve space
%{ \begingroup \vskip 10pt \begin{center} \large {\bf \@title}
%	\vskip 10pt \large \@author \hskip 20pt \@date \end{center}
%  \vskip 10pt \endgroup \setcounter{footnote}{0} }
%\makeatother % End of region containing @ commands
\renewcommand{\labelenumi}{(\alph{enumi})} % Use letters for enumerate
% \DeclareMathOperator{\Sample}{Sample}
\let\vaccent=\v % rename builtin command \v{} to \vaccent{}
\renewcommand{\v}[1]{\ensuremath{\mathbf{#1}}} % for vectors
\newcommand{\gv}[1]{\ensuremath{\mbox{\boldmath$ #1 $}}} 
% for vectors of Greek letters
\newcommand{\uv}[1]{\ensuremath{\mathbf{\hat{#1}}}} % for unit vector
%\newcommand{\abs}[1]{\left| #1 \right|} % for absolute value
\newcommand{\avg}[1]{\left< #1 \right>} % for average
%\let\underdot=\d % rename builtin command \d{} to \underdot{}
%\renewcommand{\d}[2]{\frac{d #1}{d #2}} % for derivatives
\newcommand{\fd}[2]{\frac{d #1}{d #2}} % for derivatives
\newcommand{\dd}[2]{\frac{d^2 #1}{d #2^2}} % for double derivatives
\newcommand{\pd}[2]{\frac{\partial #1}{\partial #2}} 
% for partial derivatives
\newcommand{\pdd}[2]{\frac{\partial^2 #1}{\partial #2^2}} 
\newcommand{\pddx}[3]{\frac{\partial^2 #1}{\partial #2 \partial #3}} 
% for double partial derivatives
\newcommand{\pdc}[3]{\left( \frac{\partial #1}{\partial #2}
 \right)_{#3}} % for thermodynamic partial derivatives
\newcommand{\ket}[1]{\left| #1 \right>} % for Dirac bras
\newcommand{\bra}[1]{\left< #1 \right|} % for Dirac kets
\newcommand{\braket}[2]{\left< #1 \vphantom{#2} \right|
 \left. #2 \vphantom{#1} \right>} % for Dirac brackets
\newcommand{\matrixel}[3]{\left< #1 \vphantom{#2#3} \right|
 #2 \left| #3 \vphantom{#1#2} \right>} % for Dirac matrix elements
\newcommand{\grad}[1]{\gv{\nabla} #1} % for gradient
\let\divsymb=\div % rename builtin command \div to \divsymb
\renewcommand{\div}[1]{\gv{\nabla} \cdot #1} % for divergence
\newcommand{\tdiv}[1]{\text{div} #1} % for divergence
\newcommand{\curl}[1]{\gv{\nabla} \times #1} % for curl
\newcommand{\lap}[1]{\gv{\nabla}^{2} #1} % for gradient
\let\baraccent=\= % rename builtin command \= to \baraccent
\renewcommand{\=}[1]{\stackrel{#1}{=}} % for putting numbers above =
\newtheorem{prop}{Proposition}
\newtheorem{thm}{Theorem}[section]
\newtheorem{lem}[thm]{Lemma}
\theoremstyle{definition}
\newtheorem{dfn}{Definition}
\theoremstyle{remark}
\newtheorem*{rmk}{Remark}

% ***********************************************************
% ********************** END HEADER *************************
% ***********************************************************

%% Labelled matrices, Felder and Felder style (column labels above the
%% parentheses, row labels to the left). Built on plain TeX's \bordermatrix.
%%   \lmat{col1}{col2}{row1}{row2}{a}{b}{c}{d}   a 2x2 with labels
%%   \lmatt{c1}{c2}{c3}{r1}{r2}{r3}{a}{b}{c}{d}{e}{f}{g}{h}{i} is one argument
%%   too many for \newcommand, so for 3x3 and up write \bordermatrix directly:
%%   \bordermatrix{ & c1 & c2 & c3 \cr r1 & a & b & c \cr r2 & d & e & f \cr r3 & g & h & i }
\newcommand{\lmat}[8]{\bordermatrix{ & #1 & #2 \cr #3 & #5 & #6 \cr #4 & #7 & #8 }}

\usepackage{xr-hyper}
\usepackage{booktabs}
\usepackage[colorlinks=true, linkcolor=blue, citecolor=blue, urlcolor=blue]{hyperref}
\externaldocument{AstroLagrange}
\newcommand{\Om}{\Omega}
\newcommand{\Lag}{\mathcal{L}}
\newcommand{\Ueff}{U_{\rm eff}}
\pagestyle{empty}
\setlength{\parskip}{2pt}
\setlength{\parindent}{0pt}
\setlength{\abovedisplayskip}{3pt}
\setlength{\belowdisplayskip}{3pt}
\setlength{\abovedisplayshortskip}{1pt}
\setlength{\belowdisplayshortskip}{1pt}

\begin{document}
\begin{center}
  {\large\bfseries The Lagrange points: the key equations}\\[1pt]
  {\small Equation numbers are the astronomy handout's (Section~1). Everything else is yours to fill in.}
\end{center}

\textbf{Set-up (Steps 1--3).} Two primaries $M_1$, $M_2$ a fixed distance $R$ apart, circling
their CM at angular speed $\Om$; a test mass $m \ll M_1, M_2$ moving in their plane ($Z = 0$).
With $\mu \equiv M_2/(M_1 + M_2)$:
\begin{equation*}
  r_1 = \mu R, \qquad r_2 = (1 - \mu)R \quad\text{(distances of $M_1$, $M_2$ from the CM)}, \qquad
  \Om^2 = \frac{G(M_1 + M_2)}{R^3} \quad\text{\eqref{eq:kepler3}} .
\end{equation*}

\textbf{Step 4: coordinates that rotate with the primaries.} $(X, Y)$ inertial, $(x, y)$ turned
through $\Om t$; the primaries then sit at $(-r_1, 0)$ and $(r_2, 0)$ for all time, so $U$ loses
its $t$:
\begin{gather*}
  X = x\cos\Om t - y\sin\Om t, \qquad Y = x\sin\Om t + y\cos\Om t \quad\text{\eqref{eq:rotation}}, \\
  U(x, y) = -\frac{GmM_1}{\sqrt{(x + r_1)^2 + y^2}} - \frac{GmM_2}{\sqrt{(x - r_2)^2 + y^2}} \quad\text{\eqref{eq:U2d}}, \\
  \dot X^2 + \dot Y^2 = \dot x^2 + \dot y^2 + 2\Om(x\dot y - y\dot x) + \Om^2(x^2 + y^2), \\
  \boxed{\;\Lag = \tfrac12 m(\dot x^2 + \dot y^2) + m\Om\,(x\dot y - y\dot x)
  + \tfrac12 m\Om^2(x^2 + y^2) - U(x, y)\;} \quad\text{\eqref{eq:Lrot3}: no $t$ in it.}
\end{gather*}

\textbf{Step 5: Lagrange's equations.} $\dfrac{d}{dt}\dfrac{\partial\Lag}{\partial\dot x} = \dfrac{\partial\Lag}{\partial x}$ and the same for $y$ give
\begin{equation*}
  m\ddot x = -\pd{U}{x} + m\Om^2 x + 2m\Om\dot y, \qquad
  m\ddot y = -\pd{U}{y} + m\Om^2 y - 2m\Om\dot x \quad\text{\eqref{eq:eom3}},
\end{equation*}
Newton's second law with three forces: gravity $-\grad{U}$; centrifugal $m\Om^2(x, y)$, straight
away from the CM; Coriolis $2m\Om(\dot y, -\dot x)$, perpendicular to the velocity, no work done.

\textbf{Step 6(a): where can $m$ sit still?} $\dot x = \dot y = \ddot x = \ddot y = 0$ in
\eqref{eq:eom3}: the Coriolis terms vanish with the velocities, and
\begin{equation*}
  \grad{\Ueff} = 0, \qquad \Ueff \equiv U - \tfrac12 m\Om^2(x^2 + y^2) \quad\text{\eqref{eq:equilibrium3}} .
\end{equation*}
Exactly five solutions. On the axis ($y = 0$), in units $R = 1$, $G(M_1 + M_2) = 1$, per unit mass:
\begin{equation*}
  x = \frac{(1-\mu)(x+\mu)}{|x+\mu|^3} + \frac{\mu\,(x-1+\mu)}{|x-1+\mu|^3} \quad\text{\eqref{eq:axis}}:
  \quad L_1, L_2, L_3, \text{ one in each stretch of the axis};
\end{equation*}
for small $\mu$, $L_1$ and $L_2$ sit $d \approx R\,(\mu/3)^{1/3}$ from $M_2$ \eqref{eq:hill}, the
Hill radius. Off the axis, $L_4$ and $L_5$ at the two equilateral points ($d_1 = d_2 = R$).

\textbf{Step 6(b): what is conserved?} No $t$ in $\Lag$, so
\begin{equation*}
  \Ham = \tfrac12 m(\dot x^2 + \dot y^2) + \Ueff \quad\text{\eqref{eq:jacobi3}, the Jacobi integral; not $T + U$.}
\end{equation*}

\textbf{Stability.} None of the five is a minimum of $\Ueff$ ($L_1$--$L_3$ saddles, $L_4$, $L_5$
maxima). Linearize \eqref{eq:eom3} about a point with the Coriolis term kept ($x = L_x + \xi$,
$y = L_y + \eta$; $A, B, C$ the second derivatives of $\Ueff$ there; same units as above):
\begin{gather*}
  \ddot\xi - 2\dot\eta = -A\xi - B\eta, \qquad \ddot\eta + 2\dot\xi = -B\xi - C\eta \quad\text{\eqref{eq:linear}}; \\
  \xi, \eta \propto e^{\lambda t}: \qquad
  \lambda^4 + (4 + A + C)\lambda^2 + (AC - B^2) = 0 \quad\text{\eqref{eq:char}} .
\end{gather*}
The $4$ is the Coriolis effect. A saddle ($AC - B^2 < 0$) always has a real $\lambda > 0$:
$L_1$--$L_3$ unstable for every $\mu$ (Sun--Earth $L_2$: $e$-folding time 23 days; JWST fires
thrusters every three weeks). At $L_4$, $A = -\tfrac34$, $C = -\tfrac94$, $AC - B^2 = \tfrac{27}{4}\mu(1-\mu)$:
\begin{equation*}
  \lambda^4 + \lambda^2 + \tfrac{27}{4}\mu(1-\mu) = 0 \quad\Rightarrow\quad
  \text{stable iff } 27\mu(1-\mu) < 1, \;\; \mu < 0.0385 \quad\text{\eqref{eq:gascheaucond}},
\end{equation*}
which covers Sun--Jupiter (the Trojans; libration period 147 yr) and Earth--Moon.

\begin{center}
\small
\begin{tabular}{lrrrr}
\toprule
 & $\mu$ & $L_1$ from $M_2$ & $L_2$ from $M_2$ & $R(\mu/3)^{1/3}$ \\
\midrule
Sun--Earth   & $3.00\times 10^{-6}$ & $1.49\times 10^6$ km & $1.50\times 10^6$ km & $1.50\times 10^6$ km \\
Sun--Jupiter & $9.54\times 10^{-4}$ & $5.19\times 10^7$ km & $5.43\times 10^7$ km & $5.31\times 10^7$ km \\
Earth--Moon  & $1.22\times 10^{-2}$ & $5.80\times 10^4$ km & $6.45\times 10^4$ km & $6.13\times 10^4$ km \\
\bottomrule
\end{tabular}
\end{center}
\end{document}
